\paragrafos{Φυσικοί αριθμοί, διάταξη, στρογγυλοποίηση: η σειρά που δεν τελειώνει}{Α.1.1, «Φυσικοί αριθμοί – Διάταξη φυσικών – Στρογγυλοποίηση»}

\begin{anagke}
Πρώτη μέρα στο γυμνάσιο. Στην αίθουσα της Α3 η Ζωή, ο Ηλίας, η Μυρτώ, ο Στέλιος, η Φανή και ο Βασίλης, ο πρώην πρόεδρος του μαθητικού συμβουλίου της ΣΤ’, βρίσκουν στα θρανία δύο γνωστά πρόσωπα από το άλλο δημοτικό της περιοχής: τον Νίκο και την Αγγελική, που τους έχουν δει σε αγώνες μπάσκετ και σε γιορτές. Στο τελευταίο θρανίο κάθεται ένα κορίτσι που δεν το ξέρει κανείς. Την λένε Πολυξένη, ήρθε το καλοκαίρι από τη Θεσσαλονίκη και δεν έχει πει ακόμα ούτε λέξη.

Στην είσοδο ένα χαρτί λέει: «Το σχολείο μας έχει 387 μαθητές». Ο Νίκος το φωτογραφίζει και στέλνει στη μαμά του: «Το σχολείο έχει 400 παιδιά!». Η μαμά του απαντάει με ένα χαμογελαστό «στρογγύλεψες».

Η κυρία Αλεξάνδρα, η καθηγήτρια των μαθηματικών, δεν ανοίγει βιβλίο. Γράφει στον πίνακα δύο αριθμούς: τους μαθητές του σχολείου, 387, και το υψόμετρο του Ολύμπου, 2.917 μέτρα. «Έχετε να τους πείτε σε έναν φίλο σας. Ποιον θα πείτε ακριβώς και ποιον περίπου;»

«Όλους ακριβώς», λέει ο Στέλιος. «Τι νόημα έχει το “περίπου”;» Η Ζωή σηκώνει το χέρι. «Αν πεις στον φίλο σου ότι περπάτησες 8.412 βήματα, θα θυμάται μόνο το 8 και θα μπερδεύεται με τα υπόλοιπα. Το “περίπου 8.000” το κρατάς». Ο Βασίλης χαμογελάει: «Εγώ λέω πάντα “γύρω στα τετρακόσια”. Με το μάτι». Τότε ακούγεται για πρώτη φορά η φωνή της Πολυξένης, χαμηλά: «Και πώς ξέρουμε ότι όλοι θα βγάλουμε το ίδιο “περίπου”; Αν ο Νίκος κάνει το 387 τετρακόσια, και ο Βασίλης “γύρω στα τριακόσια ογδόντα”;»

Όλοι γυρίζουν να την κοιτάξουν. Η κυρία Αλεξάνδρα χαμογελάει: η ερώτηση είναι ακριβώς αυτή που ήθελε να ακούσει. Πότε επιτρέπεται το «περίπου», και πώς γίνεται με κανόνα, ώστε όλοι να κάνουμε το ίδιο;
\end{anagke}

\begin{lexi}{φυσικός, στρογγυλοποίηση}{φυσικός, στρογγυλοποίηση}
Οι \lexiword{φυσικοί} αριθμοί λέγονται έτσι επειδή είναι οι πρώτοι που βρίσκει κανείς στη φύση, όταν μετράει πράγματα: μήλα, παιδιά, βήματα. Κανείς δεν μέτρησε ποτέ «ενάμισι παιδί». Το μηδέν μπήκε στην παρέα αργότερα. Στο σχολείο θα το βρεις μέσα στους φυσικούς.

Η \lexiword{στρογγυλοποίηση} είναι το να κάνεις έναν αριθμό «στρογγυλό», δηλαδή να τον αλλάξεις με έναν κοντινό του που τελειώνει σε μηδενικά: 387 γίνεται 400. Στρογγυλός σημαίνει τακτοποιημένος, χωρίς αιχμές. Δεν είναι μαγικό: είναι \emph{συμφωνία}. Όλοι κάνουμε τα ίδια βήματα και βγάζουμε το ίδιο αποτέλεσμα.

Δύο ακόμα λέξεις που θα ξανασυναντήσεις: \lexiword{άρτιος} είναι ο αριθμός που «ταιριάζει» τέλεια σε ζευγάρια, και \lexiword{περιττός} αυτός που «περισσεύει» ένας. Τους λέμε και ζυγούς και μονούς.
\end{lexi}

\begin{metafrasi}
Ο αριθμός που ακολουθεί το 9. & $9+1=10$\\
Το πέντε είναι μεγαλύτερο από το τρία. & $5>3$\\
Ο αριθμός $x$ βρίσκεται ανάμεσα στο 3 και στο 8. & $3<x<8$\\
Τρία εκατομμύρια, σαράντα χιλιάδες και επτά. & 3.040.007\\
Το 3.486.917 στρογγυλοποιημένο στις χιλιάδες. & $3.486.917\approx 3.487.000$\\
Ο αριθμός είναι άρτιος. & Διαιρείται με το 2\\
\end{metafrasi}

\tmima{Η θεωρία}

\begin{orismos}{φυσικοί αριθμοί}
Οι αριθμοί $0,\,1,\,2,\,3,\,4,\,\dots,\,99,\,100,\,\dots$ λέγονται \textbf{φυσικοί αριθμοί}. Κάθε φυσικός έχει έναν επόμενο και, εκτός από το 0, έναν προηγούμενο.
\end{orismos}

Οι φυσικοί χωρίζονται σε δύο κατηγορίες: \textbf{άρτιους} (ή ζυγούς), που διαιρούνται με το 2, και \textbf{περιττούς} (ή μονούς), που δεν διαιρούνται. Ο άρτιος τελειώνει πάντα σε 0, 2, 4, 6 ή 8.

\tmimasmall{Δέκα ψηφία, άπειροι αριθμοί}

Με τα δέκα ψηφία $0,\,1,\,2,\,\dots,\,9$ γράφουμε όσους αριθμούς θέλουμε. Αυτό γίνεται επειδή η αξία ενός ψηφίου εξαρτάται και από τη \textbf{θέση} του, δηλαδή από τη \textbf{δεκαδική τάξη}: μονάδες, δεκάδες, εκατοντάδες, χιλιάδες, δεκάδες χιλιάδων, εκατοντάδες χιλιάδων, εκατομμύρια. Το 0 κρατάει μια θέση άδεια: στο 6.040.205 τα τρία μηδενικά λένε ότι δεν υπάρχουν εκατοντάδες χιλιάδων, χιλιάδες και δεκάδες.

\tmimasmall{Διάταξη: ποιος είναι μεγαλύτερος}

Γράφουμε $=$ για «ίσος με», $<$ για «μικρότερος από» και $>$ για «μεγαλύτερος από». Δύο φυσικούς μπορούμε πάντα να τους συγκρίνουμε.

\begin{kanonas}{πώς συγκρίνω δύο φυσικούς}
\begin{enumerate}[leftmargin=1.6em,itemsep=0.2ex,topsep=2pt]
\item Μετράω τα ψηφία. Ο αριθμός με περισσότερα ψηφία είναι μεγαλύτερος (αν δεν ξεκινάει με 0).
\item Αν έχουν ίδιο πλήθος ψηφίων, τα κοιτάζω από αριστερά προς τα δεξιά. Στο πρώτο ψηφίο που διαφέρουν, μεγαλύτερος είναι αυτός με το μεγαλύτερο ψηφίο.
\end{enumerate}
\end{kanonas}

Έτσι $73.409>73.094$, γιατί τα δύο πρώτα ψηφία ταιριάζουν και στο τρίτο το 4 είναι μεγαλύτερο από το 0. Και $812<10.250$, γιατί το 812 έχει λιγότερα ψηφία. Τους φυσικούς μπορούμε να τους βάλουμε σε σειρά: \textbf{αύξουσα}, από τον μικρότερο στον μεγαλύτερο, ή \textbf{φθίνουσα}, αντίστροφα. Δεν υπάρχει μεγαλύτερος φυσικός αριθμός, γιατί κάθε φυσικός έχει επόμενο. Ο μικρότερος είναι το 0.

\tmimasmall{Η αριθμητική ευθεία}

Η διάταξη μας επιτρέπει να βάλουμε τους φυσικούς πάνω σε μια ευθεία. Διαλέγουμε ένα σημείο Ο για το 0 (την αρχή) και ένα σημείο Α δεξιά του για το 1. Το τμήμα ΟΑ είναι η \textbf{μονάδα μέτρησης}. Με το ίδιο μήκος βρίσκουμε τα σημεία των 2, 3, 4, \dots

\begin{center}
\begin{tikzpicture}[x=1.15cm,y=1cm]
\draw[-{Stealth}] (-0.3,0)--(9.6,0);
\foreach \i in {0,...,8}{\draw (\i,-0.12)--(\i,0.12); \node[below,font=\small] at (\i,-0.12) {\i};}
\node[above,font=\small\bfseries,gblue] at (0,0.15) {Ο};
\node[above,font=\small\bfseries,gblue] at (1,0.15) {Α};
\draw[<->,gblue] (0,0.75)--(1,0.75) node[midway,above,font=\scriptsize] {μονάδα};
\end{tikzpicture}
\end{center}

\tmimasmall{Στρογγυλοποίηση}

\begin{kanonas}{στρογγυλοποίηση φυσικού}
\begin{enumerate}[leftmargin=1.6em,itemsep=0.2ex,topsep=2pt]
\item Βρίσκω την τάξη στην οποία θα στρογγυλοποιήσω (δεκάδες, εκατοντάδες, χιλιάδες, \dots).
\item Κοιτάζω το ψηφίο της \textbf{αμέσως μικρότερης} τάξης.
\item Αν είναι $0,\,1,\,2,\,3$ ή $4$, το ψηφίο της τάξης μου μένει ίδιο. Αν είναι $5,\,6,\,7,\,8$ ή $9$, αυξάνεται κατά 1.
\item Όλα τα ψηφία δεξιά από την τάξη μου γίνονται 0.
\end{enumerate}
\end{kanonas}

Το αποτέλεσμα λέγεται \textbf{προσέγγιση} του αριθμού και γράφεται με το σύμβολο $\approx$ (διαβάζεται «περίπου ίσο»). Για παράδειγμα, $2.917\approx 3.000$ στις χιλιάδες και $2.917\approx 2.900$ στις εκατοντάδες.

Η στρογγυλοποίηση ταιριάζει εκεί που μας ενδιαφέρει η τάξη μεγέθους, όπως το πλήθος των ανθρώπων σε μια πόλη. Δεν ταιριάζει εκεί που χρειάζεται ακρίβεια: ο αριθμός του τηλεφώνου, το ΑΦΜ, ο κωδικός ενός λουκέτου.

\begin{prosoxi}
\begin{itemize}[leftmargin=1.2em,itemsep=0.3ex,topsep=2pt]
\item Στρογγυλοποιούμε \textbf{μία φορά}, από τον αρχικό αριθμό. Το 1.449 στις εκατοντάδες είναι 1.400, γιατί το ψηφίο των δεκάδων είναι 4. Αν το κάνεις σε δύο βήματα, πρώτα στις δεκάδες (1.450) και μετά στις εκατοντάδες, βγάζεις 1.500, που είναι λάθος.
\item Όταν το ψηφίο που αυξάνεται είναι 9, η αύξηση «κουβαλάει»: $9.995\approx 10.000$ στις δεκάδες.
\item Όταν μετράς \textbf{από ως}, και τα δύο άκρα μετράνε. Από τη σελίδα 20 ως τη σελίδα 47 είναι $47-20+1=28$ σελίδες, όχι 27.
\end{itemize}
\end{prosoxi}

\tmima{Λυμένα παραδείγματα}

\paradeigma{Το ντουλάπι του Ηλία}
\tixero Το ντουλάπι του Ηλία ανοίγει με τριψήφιο κωδικό. Θυμάται ότι έχει τα ψηφία 3, 0 και 5, καθένα μία φορά, αλλά όχι με ποια σειρά.\par
\tithelo Όλους τους πιθανούς κωδικούς, σε αύξουσα σειρά.\par
\possindeo Δοκιμάζω όλες τις σειρές των τριών ψηφίων: 305, 350, 035, 053, 503, 530. Ο 035 και ο 053 δεν είναι τριψήφιοι αριθμοί, γιατί το 0 μπροστά δεν μετράει. Τους σβήνω.\par
\apantisi 305, 350, 503, 530.\par
\elegxos Είναι τέσσερις, και κανένας δεν ξεκινάει από 0. Ο μικρότερος είναι ο 305: ξεκινάει από το μικρότερο μη μηδενικό ψηφίο, το 3, και μετά βάζει το μικρότερο που μένει, το 0.

\paradeigma{Οι προβολές του Βασίλη}
\tixero Το βίντεο της σχολικής εκλογικής εκστρατείας του Βασίλη έχει 3.486.917 προβολές. (Ναι, ο Βασίλης υπερβάλλει, αλλά ας το δεχτούμε.)\par
\tithelo Να το στρογγυλοποιήσω (α) στις εκατοντάδες, (β) στις χιλιάδες, (γ) στις εκατοντάδες χιλιάδων.\par
\possindeo
(α) Τάξη: εκατοντάδες. Η αμέσως μικρότερη τάξη είναι οι δεκάδες, με ψηφίο 1. Επειδή $1<5$, το 9 μένει ίδιο και τα δεξιά γίνονται 0: $3.486.917\approx 3.486.900$.\par
(β) Τάξη: χιλιάδες, ψηφίο 6. Η εκατοντάδα έχει ψηφίο 9. Επειδή $9\ge 5$, το 6 γίνεται 7: $3.486.917\approx 3.487.000$.\par
(γ) Τάξη: εκατοντάδες χιλιάδων, ψηφίο 4. Οι δεκάδες χιλιάδων έχουν ψηφίο 8. Επειδή $8\ge 5$, το 4 γίνεται 5: $3.486.917\approx 3.500.000$.\par
\apantisi 3.486.900, 3.487.000 και 3.500.000.\par
\elegxos Στο (β) το 3.486.917 απέχει 83 από το 3.487.000 και 917 από το 3.486.000. Το πρώτο είναι πιο κοντά. Στο (γ) απέχει 13.083 από το 3.500.000 και 86.917 από το 3.400.000. Και πάλι το 3.500.000 είναι πιο κοντά.

\paradeigma{Πόσες σελίδες;}
\tixero Η Φανή διαβάζει από τη σελίδα 58 ως τη σελίδα 131 του βιβλίου της, και τις δύο μαζί.\par
\tithelo Πόσες σελίδες διαβάζει.\par
\possindeo Αν έβγαζα $131-58=73$, θα είχα αφήσει έξω την πρώτη σελίδα. Για να μετρήσω και τα δύο άκρα, προσθέτω 1: $131-58+1=74$.\par
\apantisi 74 σελίδες.\par
\elegxos Δοκιμή με μικρούς αριθμούς: από τη σελίδα 3 ως τη 5 είναι τρεις σελίδες (3, 4, 5). Ο τύπος δίνει $5-3+1=3$. Σωστά.

\tmima{Ασκήσεις}

\zoni{Το κατάλαβα}{εφαρμογή του κανόνα}

\ex Γράψε με ψηφία: α) τριακόσια δεκαοκτώ \quad β) εννιακόσια ένα \quad γ) δεκαπέντε χιλιάδες διακόσια τέσσερα \quad δ) δύο εκατομμύρια τριακόσιες χιλιάδες εβδομήντα.

\ex Γράψε με λέξεις: α) 64.307 \quad β) 1.050.009 \quad γ) 200.020.

\ex Ποιοι είναι οι τρεις προηγούμενοι αριθμοί του 1.000 και ποιοι οι δύο επόμενοι του 999.999; Ποιος είναι ο προηγούμενος του 0;

\ex Τα σκορ από το παιχνίδι του Βασίλη είναι 7.250, 7.205, 72.050, 7.520 και 5.270. Γράψε τα σε αύξουσα σειρά και μετά σε φθίνουσα.

\ex Βάλε το σωστό σύμβολο ($<$, $=$ ή $>$): α) $603\ \square\ 630$ \quad β) $9.099\ \square\ 9.100$ \quad γ) $45.012\ \square\ 45.021$ \quad δ) $100.000\ \square\ 99.999$ \quad ε) $7.008\ \square\ 7.008$ \quad στ) $12.345\ \square\ 1.234$.

\ex Σχεδίασε μια αριθμητική ευθεία με αρχή το Ο και μονάδα $\mathrm{ΟΑ}=3$ εκ. Τοποθέτησε τα σημεία Β, Γ, Δ και Ε σε αποστάσεις 9, 15, 21 και 27 εκ. από το Ο. Ποιοι αριθμοί αντιστοιχούν σε αυτά;

\ex Στρογγυλοποίησε στην πλησιέστερη εκατοντάδα: 248, 350, 1.249, 5.951, 19.999 και 70.049.

\ex Στρογγυλοποίησε τον αριθμό 2.846.593 (α) στις δεκάδες, (β) στις εκατοντάδες, (γ) στις χιλιάδες, (δ) στις δεκάδες χιλιάδων, (ε) στις εκατοντάδες χιλιάδων, (στ) στα εκατομμύρια.

\ex Σωστό ή λάθος;\par
α) Στον αριθμό 4.050.300 τα μηδενικά δείχνουν ότι δεν υπάρχουν εκατοντάδες χιλιάδων, χιλιάδες, δεκάδες και μονάδες.\quad β) Εκατό δεκάδες είναι χίλιες μονάδες.\quad γ) Από τον 40 ως τον 99 υπάρχουν 60 αριθμοί.\quad δ) Από τη σελίδα 20 ως τη σελίδα 29 είναι 9 σελίδες.\quad ε) Αν σήμερα είναι Δευτέρα, σε 15 ημέρες θα είναι πάλι Δευτέρα.\quad στ) Ανάμεσα στο 7 και στο 8 δεν υπάρχει άλλος φυσικός αριθμός.\quad ζ) Το 2.449 στρογγυλοποιημένο στις εκατοντάδες είναι 2.500.

\zoni{Το χρησιμοποιώ}{προβλήματα από τη ζωή}

\ex \textbf{Το λουκέτο της Φανής.} Το λουκέτο της ποδηλατοθήκης έχει τριψήφιο κωδικό με τα ψηφία 2, 0 και 8, καθένα μία φορά. Γράψε όλους τους δυνατούς κωδικούς από τον μικρότερο στον μεγαλύτερο. Πόσους θα δοκιμάσει η Φανή στη χειρότερη περίπτωση;

\ex \textbf{Το χωριό του παππού.} Η πινακίδα στην είσοδο του χωριού γράφει «Κάτοικοι: 1.846». Ο Νίκος λέει στους φίλους του «περίπου 2.000» και η Μυρτώ «περίπου 1.800». Σε ποια τάξη στρογγυλοποίησε ο καθένας; Ποιος είναι πιο κοντά; Αν ο δήμαρχος πρέπει να τυπώσει ψηφοδέλτια για όλους τους κατοίκους, μπορεί να χρησιμοποιήσει το «περίπου»; Γιατί;

\ex \textbf{Τα βήματα του Νίκου.} Η εφαρμογή στο κινητό του Νίκου γράφει ότι σήμερα έκανε 14.992 βήματα. Ο Νίκος λέει «Στρογγυλά, 15.000, πέτυχα τον στόχο!». Ο στόχος είναι τουλάχιστον 15.000 βήματα. Έφτασε τον στόχο; Πόσα του λείπουν; Γιατί εδώ η στρογγυλοποίηση δεν βοηθάει;

\ex \textbf{Τα πλακάκια του θείου.} Ο θείος Λευτέρης χρειάζεται 2.314 πλακάκια για τον κήπο. Πωλούνται σε κιβώτια των 100. Ο Στέλιος στρογγυλοποιεί στην πλησιέστερη εκατοντάδα και λέει «23 κιβώτια». Φτάνουν; Πόσα κιβώτια πρέπει να αγοράσει και πόσα πλακάκια θα περισσέψουν;

\ex \textbf{Το βιβλίο της Μυρτώς.} Η Μυρτώ διαβάζει από τη σελίδα 204 ως τη σελίδα 311, και τις δύο μαζί. Πόσες σελίδες διαβάζει; Αν διαβάζει 9 σελίδες την ημέρα, σε πόσες ημέρες θα τελειώσει;

\ex \textbf{Ο φράχτης.} Ο παππούς φράζει ένα κομμάτι του κήπου μήκους 36 μ. Βάζει έναν πάσσαλο κάθε 3 μ., και έναν στην αρχή και έναν στο τέλος. Πόσους πασσάλους θα χρειαστεί;

\ex \textbf{Οι κάρτες του αγώνα.} Οι κάρτες του σχολικού αγώνα είναι αριθμημένες από το 501 ως το 780. Πόσες κάρτες υπάρχουν; Πόσες από αυτές έχουν περιττό αριθμό;

\ex \textbf{Ο άξονας του Βασίλη.} Ο Βασίλης σχεδιάζει έναν άξονα για τις βαθμολογίες του παιχνιδιού του. Τα σημεία 0 και 1.000 απέχουν 4 εκ. (δες το σχήμα). Ποιοι αριθμοί αντιστοιχούν στα σημεία Α, Β, Γ και Δ;

\begin{center}
\begin{tikzpicture}[x=1.15cm]
\draw[-{Stealth}] (-0.3,0)--(9.6,0);
\draw (0,-0.12)--(0,0.12) node[above,font=\small] {0};
\draw (4,-0.12)--(4,0.12) node[above,font=\small] {1.000};
\foreach \x/\n in {1/Α,3/Β,5/Γ,7.5/Δ}{\draw[gblue,thick] (\x,-0.15)--(\x,0.15); \node[below,font=\small\bfseries,gblue] at (\x,-0.15) {\n};}
\end{tikzpicture}
\end{center}

\zoni{Το σκέφτομαι}{μια ερώτηση που πάει πιο μακριά}

\ex Πόσες φορές γράφεται το ψηφίο 3, αν γράψεις όλους τους αριθμούς από το 1 ως το 200;

\ex Ποιος είναι ο μικρότερος και ποιος ο μεγαλύτερος φυσικός αριθμός που, στρογγυλοποιημένος στις χιλιάδες, δίνει 8.000;

\ex Ποιος είναι ο μεγαλύτερος και ποιος ο μικρότερος πενταψήφιος αριθμός που έχει όλα τα ψηφία του διαφορετικά;

\ex Ο Νίκος λέει: «Θα βρω τον μεγαλύτερο φυσικό αριθμό και θα τον γράψω στον πίνακα». Η Ζωή λέει: «Δεν υπάρχει». Εξήγησε με λόγια γιατί έχει δίκιο η Ζωή.

\ex Ο Ηλίας δοκιμάζει στο κινητό: $7+5=12$, $3+9=12$. Ρωτάει την Πολυξένη: «Αν προσθέσω δύο φυσικούς, βγαίνει πάντα φυσικός. Αν αφαιρέσω;» Η Πολυξένη γράφει $3-5$. Μπορεί να υπάρχει αυτό το αποτέλεσμα στους φυσικούς; Τι νομίζεις ότι θα χρειαστούμε;

\zoni{Για το σπίτι}{πιο μακριά, με χαρτί και χάρακα}

\ex \textbf{Το θερμόμετρο της Αγγελικής.} Η Αγγελική φτιάχνει ένα χάρτινο θερμόμετρο για την τάξη. Σημειώνει το $0^\circ$ και το $40^\circ$ και η απόσταση ανάμεσά τους είναι 20 εκ. Πού θα βάλει τα σημάδια για τους $10^\circ$, $20^\circ$, $25^\circ$ και $35^\circ$;

\ex Δοκίμασε τις αφαιρέσεις $9-3$, $15-7$ και $21-5$. Τι είδους αριθμούς βρίσκεις; Εξήγησε γιατί η διαφορά δύο περιττών αριθμών είναι πάντα άρτιος.

\ex Πόσοι τριψήφιοι αριθμοί υπάρχουν; (Βοήθεια: ο μικρότερος είναι ο 100 και ο μεγαλύτερος ο 999, και τα δύο άκρα μετράνε.)
